\documentclass[../main.tex]{subfiles}

\begin{document}
    \subsection{{\Statement} {\Control}}
    \subsubsection{Percabangan \textit{If}}
    \paragraph{\textit{If} Tunggal: mengerjakan perintah pada cabang program hanya jika kondisi terpenuhi.}
\begin{indentlevel1}
\begin{lstlisting}
 if (kondisi) { statemen }
\end{lstlisting}
\end{indentlevel1}
    \paragraph{\textit{If-else}: mengerjakan perintah pada kondisi \textit{if} bernilai benar atau \textit{if} bernilai salah yang ditetapkan sebagai \textit{else}.}
\begin{indentlevel1}
\begin{lstlisting}
 if (kondisi) { statemen }
 else (kondisi) { statemen }
\end{lstlisting}
\end{indentlevel1}
    \paragraph{\textit{Nested if}: percabangan pada \textit{if-else} yang memiliki lebih dari 2 pilihan.}
\begin{indentlevel1}
\begin{lstlisting}
 if (kondisi) { statemen }
 else if (kondisi) { statemen }
 else (kondisi) { statemen }
\end{lstlisting}
\end{indentlevel1}
    \subsubsection{\textit{Switch case}: merupakan bentuk lain dari \textit{nested if}, hal ini karena \textit{switch case} memiliki prinsip kerja yang sama, tetapi dengan format penulisan yang berbeda \parencite{hanief2020konsep}.}
\begin{indentlevel0}
\begin{lstlisting}
 switch (variabel) {
    case 0:
        statemen
        break;
    case 1:
        statemen
        break;
    ...
    case n:
        statemen
        break;
    default:
        statemen
        break;
 }
\end{lstlisting}
\end{indentlevel0}

    \subsection{{\Looping}}
    Struktur perulangan dapat dibedakan menjadi tiga jenis yaitu \textit{for}, \textit{while} dan \textit{dowhile} \parencite{anam2021cara}.

    \subsubsection{\textit{For}: perulangan yang batasan dan syarat pengulangannya sudah ditentukan \parencite{suryadi2018belajar}.}
\begin{indentlevel0}
\begin{lstlisting}
 for (inisialisasi; kondisi perulangan; pengubah nilai cacah)
\end{lstlisting}
\end{indentlevel0}
    \subsubsection{\textit{While}: perulangan yang akan dijalankan selama kondisi terpenuhi dengan pemeriksaan kondisi sebelum dijalankan \parencite{anam2021cara}.}
\begin{indentlevel0}
\begin{lstlisting}
 while (kondisi) {
    statemen
 }
\end{lstlisting}
\end{indentlevel0}
    \subsubsection{\textit{Do-while}: perulangan yang dijalankan terlebih dahulu 1 kali, kemudian akan ada pemeriksaan kondisi \textit{while}, jika terpenuhi perulangan terjadi \parencite{anam2021cara}.}
\begin{indentlevel0}
\begin{lstlisting}
    do {
        statemen
    } while (kondisi)
\end{lstlisting}
\end{indentlevel0}
\end{document}
